% ===== latex-recomposed:document =====
% flm: style file flmrecomposed.7a0d2a74.sty
The Polish word język and the Russian word язык both mean language.
%%
% ===== latex-recomposed:style flmrecomposed.7a0d2a74.sty =====
\NeedsTeXFormat{LaTeX2e}
\ProvidesPackage{flmrecomposed.7a0d2a74}
% --- packages ---
% Some lines below differ by TeX engine. They are written inside tests of the
% engine, which the package iftex defines: \iftutex is true under LuaLaTeX and
% XeLaTeX, \ifluatex under LuaLaTeX, and \ifxetex under XeLaTeX. The last
% branch of a test is for pdfLaTeX and any other engine.
\RequirePackage{iftex}
% flm.latex.pkg.newunicodechar
\RequirePackage{newunicodechar}
\iftutex
% flm.latex.pkg.untechxt.fontenc-t2a
\RequirePackage[T2A,TU]{fontenc}
\else
% flm.latex.pkg.untechxt.fontenc-t1, flm.latex.pkg.untechxt.fontenc-t2a
\RequirePackage[T1,T2A]{fontenc}
\fi
% --- the default font encoding ---
% fontenc makes the last encoding of its option list the default, but the
% options above may be merged from several requests, so their order carries
% no meaning. The line below therefore sets the default for each engine whose
% own encoding the request loads: T1 under pdfLaTeX, TU under LuaLaTeX and
% XeLaTeX. A letter of another encoding switches to that encoding by itself.
\iftutex
\renewcommand{\encodingdefault}{TU}
\else
\renewcommand{\encodingdefault}{T1}
\fi
% --- unicode characters written as UTF-8 ---
\newunicodechar{ę}{\k{e}}
\newunicodechar{з}{{\fontencoding{T2A}\selectfont\cyrz}}
\newunicodechar{к}{{\fontencoding{T2A}\selectfont\cyrk}}
\newunicodechar{ы}{{\fontencoding{T2A}\selectfont\cyrery}}
\newunicodechar{я}{{\fontencoding{T2A}\selectfont\cyrya}}
\endinput
